\section{A second signal: persistent volatility richness}\label{sec:5}

\subsection*{Why sell the names where realized volatility keeps disappointing?}

Some names repeatedly realize less movement than their option premiums imply. The second signal uses that persistence within each name rather than requiring a second company with offsetting vega. It asks whether an observed premium has been sufficiently persistent to justify selling the next roughly 30-day volatility exposure. The economic idea is related to Goyal and Saretto \cite{goyal2009}; the new question here is how to size and allocate around it.


For every historical candidate, subtract entry ATM-call IV from the annualized RV measured after the next-session entry close through that contract's expiry. Only after the outcome finishes can it join the signal history. Average those completed RV - IV observations over the preceding 365 calendar days, requiring 126. Rank eligible names each day and sell the bottom 10\%.


\begin{equation*}M_{i,t}=\operatorname{mean}_{u:\,t-365<T_u<t}\left(RV_{i,u\to T_u}-IV^{\rm ATM\ call}_{i,u}\right).\end{equation*}

For illustration, a name whose completed outcomes average RV 26\% against entry IV 34\% has memory -8 volatility points. A name averaging -2 points ranks less rich under this rule, although neither figure guarantees a profitable new short. Ranking is cross-sectional: bottom decile means the lowest eligible memories, not a fixed requirement that the premium exceed eight points.


\subsection*{Why a straddle, and what risk does it sell?}

ATM means the strike is near the current stock price. We sell an ATM call and put with the same name, strike and expiry, then adjust the stock delta hedge daily. This collects premium against movement in either direction. At full size, \$450 is the combined call-and-put vega budget. The position is short gamma and vega: quiet realized paths can support carry, while jumps and changing uncertainty can generate large losses. It is not vega neutral.


\begin{table}[H]\centering\small
\caption{Economic exposures of the two retained signals.}
\setlength{\tabcolsep}{4pt}\renewcommand{\arraystretch}{1.18}
\begin{tabularx}{\textwidth}{@{}>{\hsize=0.80859\hsize\linewidth=\hsize\raggedright\arraybackslash}X>{\hsize=1.09570\hsize\linewidth=\hsize\raggedright\arraybackslash}X>{\hsize=1.09570\hsize\linewidth=\hsize\raggedright\arraybackslash}X@{}}\toprule
\textbf{Economic distinction} & \textbf{Relative put pair} & \textbf{ATM straddle}\\\midrule
Comparison & Current gap versus pair history & Completed RV - IV versus other names\\
Exposure & Buy one put; sell another & Sell call and put on one name\\
Initial vega & Approximately offset across names & Net short\\
Main challenge & Different realized paths and events & Premium income versus movement risk\\
\bottomrule\end{tabularx}\end{table}

The signal retains ATM-call IV even though execution uses a straddle. It ranks a difference in annualized volatility units, not squared variance and not an exact option-return forecast. The actual-contract ledger determines profitability after premiums, hedging, costs and financing. This distinction prevents the ranking target from being mistaken for simulated profit.


The two signals consequently offer different portfolios. The pair trade targets relative option richness with initial vega offset; the straddle accepts market-volatility exposure to collect a persistent premium. Their different dollar profits and risk make capital allocation a separate decision from either trade's entry size.
